\label{sec:conclusao}

\chapter{Conclusão}


Apesar do avanço das tecnologias \ac{open_source} nas últimas décadas, muitas instituições ainda apresentam uma certa resistência em incorporar essas soluções em lugar das ferramentas proprietárias. Fatores relacionados à assistência técnica, à confiabilidade, à segurança, à continuidade, manutenção e à complexidade de implementação muitas vezes levam a gestores a escolher plataformas comerciais, mesmo diante dos altos custos de licenciamento.

A análise da literatura efetuada neste estudo evidenciou que o ecossistema \ac{open_source} conta, atualmente, com uma grande quantidade de ferramentas para todas as fases de um pipeline de dados. Tecnologias como Apache Kafka, Apache Airflow, Apache Spark, Apache NiFi, HDFS e Apache Superset são amplamente empregadas, proporcionando funcionalidades equivalentes às disponíveis em várias soluções proprietárias. Entretanto, as pesquisas também demonstram que a implementação dessas tecnologias não se configura como uma tarefa simples, demandando conhecimento técnico, integração entre diversos componentes, definição de arquiteturas e um esforço adicional para configuração e manutenção. Assim, a grande quantidade de ferramentas acessíveis pode constituir um obstáculo para organizações que almejam implementar iniciativas de ED com base unicamente em tecnologias \ac{open_source}.

Visando auxiliar na diminuição dessa barreira, esta pesquisa apresentou uma proposta de arquitetura de pipeline de dados fundamentada em ferramentas \ac{open_source}, elaborada com base nos resultados da revisão da literatura e concretizada através de um protótipo funcional. A arquitetura sugerida evidenciou a possibilidade de combinar diversas tecnologias \ac{open_source} para criar uma solução integrada e eficiente, capaz de atender às necessidades de um processo de dados completo.


\section{Ameaças à Validade}

Fazendo uma análise crítica dos resultados e conclusões deste trabalho, é importante reconhecer as principais ameaças à validade que podem impactar a interpretação dos achados apresentados. Essas ameaças estão relacionadas a diferentes aspectos do processo de pesquisa e desenvolvimento, incluindo a forma como os dados foram coletados, analisados e interpretados, bem como a generalização dos resultados para outros contextos \cite[]{ihantola2011threats}.

\begin{itemize}

\item \textbf{Validade de Conclusão:}  As conclusões deste trabalho foram fundamentadas principalmente na revisão da literatura e na implementação de um protótipo funcional. Não sendo realizados experimentos controlados, \textit{benchmarks} de desempenho ou estudos comparativos que permitissem mensurar quantitativamente aspectos como escalabilidade, confiabilidade, custo operacional ou desempenho da arquitetura proposta. Dessa forma, as conclusões devem ser interpretadas como evidências de viabilidade técnica e não como comprovação definitiva.


\item \textbf{Validade Interna :} O protótipo foi desenvolvido em um ambiente controlado e utilizando um conjunto específico de fontes de dados relacionadas ao setor de turismo. A ausência de comparação direta com arquiteturas proprietárias ou outras arquiteturas open source limita a capacidade de atribuir os resultados exclusivamente às tecnologias selecionadas. Além disso, fatores como configuração de hardware, infraestrutura disponível e características particulares das fontes de dados podem influenciar os resultados obtidos durante a implementação.

\item \textbf{Validade de Construção:} A arquitetura proposta foi construída com base nos resultados da revisão rápida da literatura e em um cenário motivador de coleta de dados turísticos. Entretanto, o protótipo não contempla todos os desafios encontrados em ambientes produtivos, tais como processamento de grandes volumes de dados, alta disponibilidade, recuperação de desastres, monitoramento avançado, segurança corporativa e governança de dados. Dessa forma, existe o risco de que algumas simplificações adotadas durante o desenvolvimento não representem integralmente a complexidade encontrada em ambientes reais de produção.

\item \textbf{Validade Externa:}   O cenário utilizado concentrou-se na coleta de dados de websites de turismo internacional, podendo apresentar características distintas de outros domínios de aplicação da ED.
Embora as ferramentas selecionadas sejam amplamente utilizadas no ecossistema de \textit{Big Data}, diferentes organizações podem possuir requisitos específicos relacionados a volume de dados, infraestrutura, governança, requisitos regulatórios e restrições operacionais.
\end{itemize}

Por fim, destaca-se que a realização de uma avaliação experimental formal não fez parte do escopo deste Trabalho de Conclusão de Curso. O objetivo principal foi identificar, por meio da literatura, ferramentas open source relevantes para pipelines de dados, propor uma arquitetura baseada nesses achados e demonstrar sua viabilidade por meio da implementação de um protótipo funcional. Como trabalhos futuros, recomenda-se a realização de benchmarks de desempenho, comparações com soluções proprietárias e estudos de caso em ambientes reais de produção.
